\begin{exercice}\label{exo1-2} (\minsyndical)

Une entreprise qui fabrique des %proth\`eses de hanche 
dispositifs \'electroniques de mesure de la tension art\'erielle
d\'ecide de faire des tests de fiabilit\'e d'un nouveau mod\`ele d\'evelopp\'e. 
Un \'echantillon de 10 %proth\`eses 
appareils est pr\'elev\'e. % et %chacune d'entre elles est soumise \`a des contraintes m\'ecaniques % r\'ep\'et\'ees.
On note le temps (en journ\'ees) 
au bout duquel on observe un dysfonctionnement sur chaque appareil. 
%(en heures) au bout duquel la proth\`ese casse. 
On obtient les valeurs suivantes:

\[
91.6 \quad 35.7 \quad 251.3 \quad 24.3 \quad 5.4 \quad 67.3 \quad 170.9 \quad 9.5 
\quad 118.4 \quad 57.1
\]

\begin{enumerate}

\item On veut repr\'esenter les observations \`a l'aide d'un histogramme \`a pas fixe. Utiliser la r\`egle de Sturges pour d\'eterminer le nombre de classes et choisir avec soin la valeur des bornes extr\^emes $a_0$ et $a_k$. Compl\'eter ensuite le tableau qui suit :\\

\begin{tabular}{|c|cc}
\hline
Classe $j$ & $] a_0 , a_1 ]$ & \ldots \\ \hline
Effectifs $n_j$ & \ldots & \ldots \\ \hline
Fr\'equences $n_j/n$ & \ldots & \ldots \\ \hline
Hauteurs $n_j/(nl)$ & \ldots & \ldots \\ \hline 
\end{tabular}\\

\item Tracer ensuite l'histogramme \`a pas fixe, et le polygone des fr\'equences associ\'e.

\item On veut ensuite tracer un histogramme \`a pas variable, de telle sorte \`a ce que chaque
classe aie le m\^eme effectif : 2 observations. Compl\'eter en cons\'equence le tableau qui suit :\\

\begin{tabular}{|c|cc}
\hline
Classe $j$ & $] a_0 , a_1 ]$ & \ldots \\ \hline
Largeurs $l_j$ & \ldots & \ldots \\ \hline 
Effectifs $n_j$ & \ldots & \ldots \\ \hline
Fr\'equences $n_j/n$ & \ldots & \ldots \\ \hline
Hauteurs $n_j/(n\,l_j)$ & \ldots & \ldots \\ \hline 
\end{tabular}\\

\item Tracer l'histogramme \`a pas variable demand\'e, et le polygone des fr\'equences associ\'e.

\item Quels sont les avantages et les inconv\'enients des deux types de repr\'esentation ?

\end{enumerate}

%\corrref{TD1_2}

\end{exercice}
